\unit{Exposición 182 — Geometría formal y geometría algebraica}
\sid{182}
\src{193}
\seminar{(Mayo de 1959)}
\paperhead{GEOMETRÍA FORMAL Y GEOMETRÍA ALGEBRAICA}%
  {por Alexander GROTHENDIECK}

\fsection{1}{Esquemas}

Se sabe que un espacio algebraico afín definido sobre un cuerpo $k$ está
esencialmente determinado por su álgebra afín $A$ (anillo de las funciones
regulares definidas sobre $k$), pues los morfismos de espacios algebraicos
$X\longrightarrow Y$ corresponden biunívocamente a los homomorfismos de
$k$-álgebras $A(Y)\longrightarrow A(X)$. El álgebra afín correspondiente a un
espacio algebraico es una $k$-álgebra de tipo finito y, desde el punto de vista
«clásico», no tiene elementos nilpotentes; recíprocamente, toda álgebra de este
tipo se obtiene como álgebra afín de un espacio algebraico definido sobre $k$.
Se dispone así de un conocido diccionario que permite interpretar en términos
de álgebra conmutativa las situaciones relativas a espacios algebraicos afines.
Desde hace mucho se ha observado que de este modo se obtienen enunciados más
generales, pues por lo común ya no es necesario suponer que los anillos en juego
sean del tipo recién considerado, y la hipótesis noetheriana suele bastar. En
particular, haya o no un cuerpo de base dado, no hay razón para excluir el caso
en que estos anillos contengan elementos nilpotentes. Hasta ahora, los geómetras
se habían negado a tener en cuenta estas indicaciones y se han obstinado en
limitarse a considerar álgebras afines sin elementos nilpotentes, es decir,
espacios algebraicos cuyos haces estructurales no tienen elementos nilpotentes
(y, las más de las veces, incluso espacios algebraicos «absolutamente
irreducibles»). El conferenciante considera que esta actitud ha constituido un
serio obstáculo para el desarrollo de métodos verdaderamente naturales en
Geometría algebraica.

Sea $A$ un anillo conmutativo. Es bien sabido que el conjunto
$X=\operatorname{Spec}(A)$ de los ideales primos de $A$ está provisto de una
topología natural, la «topología de Zariski» o topología espectral. Por otra
parte, hay un haz de anillos conmutativos $\Osh_X$ sobre $X$, cuyo tallo en
$\mathfrak p\in X$ es el anillo localizado $A_{\mathfrak p}$ y cuyo anillo de
secciones se identifica con $A$. Así, $X$ se convierte en un \emph{espacio
anillado}, llamado \emph{espectro primo} de $A$. Un homomorfismo de anillos
$f:A\longrightarrow B$ define un morfismo de espacios anillados
$f':\operatorname{Spec}(B)\longrightarrow\operatorname{Spec}(A)$; la
aplicación subyacente de conjuntos no es otra que
$\mathfrak p\longmapsto f^{-1}(\mathfrak p)$. Los morfismos

\src{194}
de espacios anillados de $\operatorname{Spec}(B)$ en
$\operatorname{Spec}(A)$ obtenidos de este modo son exactamente aquellos para
los cuales los homomorfismos $\Osh_x\longrightarrow\Osh_y$ ($x=f'(y)$) son
locales (es decir, la imagen inversa del ideal maximal es el ideal maximal).

Se llama \emph{esquema afín} a un espacio anillado isomorfo a algún
$\operatorname{Spec}(A)$, y \emph{preesquema} a un espacio anillado localmente
afín, es decir, tal que cada punto tiene un entorno abierto que, con la
estructura inducida, es un esquema afín. Los \emph{morfismos} de preesquemas se
definen de manera evidente; localmente corresponden a homomorfismos de anillos.

Cuando se fija un preesquema $S$ y se consideran morfismos de preesquemas
$X\longrightarrow S$, el espacio $S$ desempeña el papel de un cuerpo o de un
anillo de base (o, mejor, de un espacio de base en una fibración). Se dice
entonces que $X$ es un $S$-preesquema; si $S=\operatorname{Spec}(A)$, esto
significa también que $\Osh_X$ es un haz de $A$-álgebras. Así, todo preesquema
puede considerarse de manera única como un $\mathbf Z$-preesquema. Naturalmente,
los $S$-preesquemas forman una categoría; además, se demuestra que en esta
categoría siempre existe el producto de dos objetos $X,Y$, que se denota por
$X\times_S Y$. Esta noción de producto permite definir el cambio de base de un
$S$-preesquema correspondiente a un morfismo $S'\longrightarrow S$: en efecto,
$X\times_S S'$ puede considerarse como un $S'$-preesquema.

Se dice que $X$ es \emph{separado} sobre $S$ si la diagonal de
$X\times_S X$ es cerrada. Se llama \emph{esquema} a un preesquema separado
sobre $\mathbf Z$; es entonces separado sobre cualquier base. Para simplificar,
en adelante solo hablaremos de esquemas y, además, los supondremos
\emph{noetherianos}, es decir, uniones finitas de abiertos afines que son
espectros de anillos noetherianos. Se dice que $X$ es \emph{de tipo finito
sobre} $S$ si, para todo abierto afín $U$ de $S$, su imagen inversa en $X$ es
una unión finita de abiertos afines cuyos anillos son álgebras de tipo finito
sobre el anillo de $U$. Son estos $S$-esquemas los que se prestan a un estudio
propiamente geométrico. En particular, para todo $s\in S$, la fibra
$f^{-1}(s)$ de $X$ sobre $s$ es un esquema algebraico sobre el cuerpo residual
$\kappa(s)$ del anillo local $\Osh_s$ de $s$ en $S$. Así, en cierta medida, $X$
puede considerarse como una familia de «espacios algebraicos» $f^{-1}(s)$, donde
el parámetro $s$ recorre $S$ (es decir, desde el punto de vista local, el
conjunto de los ideales primos de un anillo dado). Naturalmente, los
$\kappa(s)$ pueden tener características distintas. Si
$S=\operatorname{Spec}(k)$, donde $k$ es un cuerpo, se recupera esencialmente la
noción usual de «espacio algebraico», con la única diferencia de que ahora el
haz estructural puede tener elementos nilpotentes.

\src{195}
Inspirándose en nociones bien conocidas, se define la noción de
\emph{morfismo proyectivo} y, más generalmente, la de \emph{morfismo propio}.
Tal morfismo es de tipo finito; además, transforma partes cerradas en partes
cerradas y conserva esta propiedad bajo cualquier cambio de base.

Siendo $X$ un esquema (noetheriano, como siempre), el haz $\Osh_X$ es un
\emph{haz coherente de anillos} en el sentido de [2]. Por tanto, los haces
coherentes de módulos sobre $X$ son también los haces que localmente son
isomorfos al conúcleo de un homomorfismo
$\Osh_X^{m}\longrightarrow\Osh_X^{n}$.

\fsection{2}{Esquemas formales}

Sean $X$ un esquema y $X'$ una parte cerrada de $X$. Entonces existe un
subhaz coherente $J$ de $\Osh_X$ tal que
$X'=\operatorname{Supp}(\Osh_X/J)$ (e incluso existe uno máximo). Provisto de
$\Osh_X/J$, el espacio $X'$ se convierte en un esquema, denotado por $X_0$;
tal esquema se llama «subesquema cerrado de $X$». Para todo $n$ también se
puede considerar $X'$ provisto de $\Osh_X/J^{n+1}$, que se denota por $X_n$;
es un subpreesquema cerrado de $X$ cuyo conjunto subyacente sigue siendo $X'$,
pero que tiene otro haz estructural, a saber,
$\Osh_{X_n}=\Osh_X/J^{n+1}$. Evidentemente, los $\Osh_{X_n}$ forman un sistema
proyectivo de haces de anillos sobre $X'$, cuyo límite proyectivo
$\overline{\Osh}_X$ se llama \emph{completación formal de $\Osh_X$ a lo largo
de $X'$}. Provisto de este haz de anillos, $X'$ se llama \emph{completación
formal de $X$ a lo largo de $X'$}; es, pues, un espacio anillado, pero en
general no un esquema. Para todo haz coherente $F$ sobre $X$ se puede considerar
análogamente la completación formal $\overline F=\varprojlim_n F_n$ de $F$ a lo
largo de $X'$ (donde $F_n=F\otimes_{\Osh_X}\Osh_X/J^{n+1}$); es un haz de
módulos sobre $\overline X$. Sus secciones se llaman «secciones formales de $F$
a lo largo de $X'$» y se identifican con los elementos de
$\varprojlim_n\Gamma(X',F_n)$. Para $F=\Osh_X$ se obtienen las «funciones
holomorfas» de $X$ a lo largo de $X'$ en el sentido de Zariski (terminología que
no seguiremos por sus interferencias con la terminología clásica).

Se llama \emph{esquema formal} (se sobreentiende: noetheriano) a un espacio
topológico $\mathfrak X$, provisto de un haz de anillos topológicos
$\Osh_{\mathfrak X}$ que satisface la condición siguiente: hay un isomorfismo
de haces de anillos topológicos
$\Osh_{\mathfrak X}=\varprojlim_n\Osh_n$, donde los $\Osh_n$ forman un sistema
proyectivo de haces de anillos sobre $\mathfrak X$, cada uno de los cuales hace
de $\mathfrak X$ un esquema $\mathfrak X_n$, y tal que, para $m\geq n$, el
homomorfismo $\Osh_m\longrightarrow\Osh_n$ es sobreyectivo y tiene por núcleo
$J_m^{n+1}$, donde $J_m$ es

\src{196}
el núcleo de $\Osh_m\longrightarrow\Osh_0$. Se demuestra entonces que
$\Osh_{\mathfrak X}$ es un \emph{haz coherente de anillos locales
noetherianos}.

En virtud de las definiciones, una completación formal $\overline X$ como la
anterior es un esquema formal y, recíprocamente, todo esquema formal es
\emph{localmente} de este tipo. De hecho, dar un esquema formal \emph{afín} (es
decir, tal que $\mathfrak X_0$ sea afín, lo cual implica que todos los
$\mathfrak X_n$ lo son) equivale a dar un anillo topológico noetheriano,
separado y completo para la topología $J$-ádica.

Las definiciones habituales de morfismo, morfismo de tipo finito, morfismo
propio, etc., para los esquemas ordinarios se extienden sin dificultad a los
esquemas formales.

\fsection{3}{Los tres teoremas fundamentales}

Sea $f:X\longrightarrow Y$ un morfismo propio de esquemas (noetherianos, como
siempre), sea $Y'$ una parte cerrada de $Y$ y sea $X'$ su imagen inversa en
$X$; consideremos las completaciones formales correspondientes $\overline Y$ y
$\overline X$. Entonces $f$ induce un morfismo
$\overline f:\overline X\longrightarrow\overline Y$ de esquemas formales, que
además es propio. Si $F$ es un haz coherente sobre $X$, entonces $\overline F$
es un haz coherente sobre $\overline X$. En el teorema 1 se olvidan $X,Y,F$ y
solo se consideran el morfismo propio $\overline f$ de esquemas formales y el
haz coherente $\overline F$ sobre $\overline X$. (No obstante, el conferenciante
solo ha escrito una demostración completa para el caso en que se parte de
$X,Y,f,F$.)

\result{TEOREMA}{1 (teorema de finitud)}
\begin{enumerate}
  \item[\textit{i.}] Los $R^{q}\overline f_{*}(\overline F)$ son haces
    coherentes sobre $\overline Y$.
  \item[\textit{ii.}] Los homomorfismos naturales
    \[
      R^{q}\overline f_{*}(\overline F)
      \longrightarrow \varprojlim_n R^{q}f_{n*}(F_n)
    \]
    son isomorfismos.
\end{enumerate}

En este enunciado se supone elegido un subhaz coherente $J$ de $\Osh_Y$ que
define $Y'$; por imagen inversa se obtiene un subhaz coherente de $\Osh_X$ que
define $X'$ y, por consiguiente, se definen $F_n,X_n,Y_n$ y
$f_n:X_n\longrightarrow Y_n$ como en el n\textsuperscript{o} 2. Si se parte de
un morfismo propio cualquiera entre dos esquemas formales, son evidentes los
pequeños cambios que deben hacerse al explicitar estas notaciones.

El teorema 1 se refería únicamente a la «cohomología formal». El teorema siguiente

\src{197}
la relaciona con la «cohomología algebraica» y es análogo a un conocido teorema
de Serre [4] sobre la comparación entre cohomología algebraica y cohomología
analítica.

\result{TEOREMA}{2 (Primer teorema de comparación)} Los
$R^{q}f_{*}(F)$ son haces coherentes sobre $Y$ (lo cual es un caso particular
del teorema 1), y los homomorfismos naturales
\[
  \overline{R^{q}f_{*}(F)}
  \longrightarrow \varprojlim_n R^{q}f_{n*}(F_n)
\]
son isomorfismos.

\result{COROLARIO}{1} Se tienen isomorfismos canónicos:
\[
  \overline{R^{q}f_{*}(F)}=R^{q}\overline f_{*}(\overline F).
\]

Para $q=0$, este corolario es una generalización del «teorema fundamental de
las funciones holomorfas» de Zariski, del que deduciremos una generalización del
«teorema de conexión» de Zariski. Obsérvese además que, mientras el teorema 1
(\textit{ii}) es trivial para $q=0$, en modo alguno ocurre lo mismo con el
teorema 2 ni con su formulación equivalente (corolario 1). En efecto, la
demostración procede por inducción descendente sobre $q$ (es trivial para $q$
grande, pues entonces ambos miembros son nulos), de modo que el caso $q=0$
aparece como el último paso de la inducción y, por así decirlo, como el caso «más
difícil».

\result{COROLARIO}{2} Supongamos $Y=\operatorname{Spec}(A)$ y que $Y'$ está
definido por el ideal $J$ de $A$. Entonces, para todo haz coherente $F$ sobre
$X$, los $H^{q}(X,F)$ son $A$-módulos de tipo finito cuyas completaciones
$J$-ádicas son los $H^{q}(\overline X,\overline F)$.

Aplicando, por último, este corolario a
$H=\Homsh_{\Osh_X}(F,G)$, se obtiene:

\result{COROLARIO}{3} Supongamos $Y=\operatorname{Spec}(A)$ y que $Y'$ está
definido por el ideal $J$ de $A$. Sean $F,G$ dos haces coherentes sobre $X$;
entonces $\operatorname{Hom}(F,G)$ es un módulo de tipo finito sobre $A$, cuya
completación $J$-ádica se identifica con
$\operatorname{Hom}(\overline F,\overline G)$.

Naturalmente, la aplicación natural
$\operatorname{Hom}(F,G)\longrightarrow
\operatorname{Hom}(\overline F,\overline G)$ es la que asocia a un
homomorfismo $u:F\longrightarrow G$ su prolongación «por continuidad»
$\overline u:\overline F\longrightarrow\overline G$ (con lo cual
$\overline F$ se convierte en un funtor de $F$).

Supongamos ahora que $A$ es separado y completo para su topología $J$-ádica;
entonces los corolarios 2 y 3 precedentes dan:
\[
  H^{q}(X,F)=H^{q}(\overline X,\overline F),
  \qquad
  \operatorname{Hom}(F,G)
  =\operatorname{Hom}(\overline F,\overline G).
\]

\src{198}
Esta última identidad muestra que la categoría de los haces coherentes sobre
$X$ se identifica con una \emph{subcategoría} (cuyos morfismos son los
morfismos inducidos) de la categoría de los haces coherentes sobre
$\overline X$. De hecho, se tiene incluso:

\result{TEOREMA}{3} Para que un haz de módulos sobre $\overline X$ sea
coherente, es necesario y suficiente que sea isomorfo a un haz de la forma
$\overline F$, donde $F$ es un haz coherente sobre $X$ (determinado salvo
isomorfismo canónico en virtud del teorema 2, corolario 3). [Recordemos que
ahora $Y=\operatorname{Spec}(A)$, siendo $A$ un anillo topológico noetheriano,
separado y completo para la topología $J$-ádica.]

\result{COROLARIO}{1} Los subesquemas cerrados de $X$ corresponden
biunívocamente a los subesquemas formales cerrados de $\overline X$.

En efecto, corresponden a los subhaces coherentes de $\Osh_X$ y,
respectivamente, de $\overline{\Osh}_X$. Considerando los grafos de morfismos
como subesquemas cerrados, se deduce del corolario 1:

\result{COROLARIO}{2} Sean $X,Z$ dos esquemas propios sobre $A$ (anillo
noetheriano, separado y completo para la topología $J$-ádica). Entonces la
aplicación $g\longmapsto\overline g$ define una correspondencia biunívoca entre
los $Y$-morfismos de $X$ en $Z$ y los $\overline Y$-morfismos de $\overline X$
en $\overline Z$.

En otros términos, los esquemas algebraicos propios sobre $A$ aparecen como una
\emph{subcategoría} (cuyos morfismos son los morfismos inducidos) de la
categoría de los esquemas formales propios sobre $\overline Y$. Sin embargo,
debe advertirse que existen esquemas formales propios sobre $\overline Y$ que
no son «algebraizables», es decir, isomorfos a algún $\overline X$ con $X$
propio sobre $A$ (del mismo modo que existen variedades analíticas complejas
compactas que no provienen de variedades algebraicas definidas sobre el cuerpo
de los complejos). Tales esquemas formales aparecen de manera natural en la
«teoría de módulos». Señalemos, no obstante, un caso particular interesante en
el que un esquema formal es algebraizable:

\result{TEOREMA}{4} Sea $A$ un anillo noetheriano local completo de cuerpo
residual $k$, y sea $\mathfrak X$ un esquema formal propio sobre $A$ (provisto
de su topología $\mathfrak m(A)$-ádica). Se supone:
\begin{enumerate}
  \item[\textit{i.}] que los anillos locales de $\Osh_{\mathfrak X}$ son
    módulos planos sobre $A$ o, lo que equivale a lo mismo, si se dota a
    $\Osh_{\mathfrak X}$ y a $A$ de la filtración definida por las potencias del
    ideal maximal de $A$, para los graduados asociados se tiene la relación
    \[
      \operatorname{gr}(\Osh_{\mathfrak X})
      \simeq
      \operatorname{gr}^{0}(\Osh_{\mathfrak X})
      \otimes_k\operatorname{gr}(A).
    \]
  \item[\textit{ii.}] Considerando
    $\mathfrak X_0=\mathfrak X\otimes_A k$ como un esquema algebraico sobre
    $k$, se tiene
    \[
      H^2(\mathfrak X_0,\Osh_{\mathfrak X_0})=0.
    \]
  \item[\textit{iii.}] $\mathfrak X_0$ es proyectivo.
\end{enumerate}
Bajo estas condiciones, $\mathfrak X$ es algebraizable y, más precisamente, es
isomorfo a $\overline X$, donde $X$ es un $A$-esquema proyectivo.

Las condiciones \textit{(ii)} y \textit{(iii)} se satisfacen, en particular,
si $\mathfrak X_0$ es una curva simple sobre $k$, y el teorema 4 puede aplicarse
en particular en la «teoría de módulos» para las curvas de género dado\ldots{}
Indiquemos cómo se demuestra el teorema 4: se demuestra (cf. la proposición 3
más adelante) que \textit{(i)} y \textit{(ii)} implican que todo haz coherente
sobre $\mathfrak X_0$, localmente isomorfo al haz fundamental, se obtiene por
reducción a partir de un haz de la misma naturaleza sobre $\mathfrak X$. Se
parte entonces de un haz «amplio» sobre $\mathfrak X_0$ (existe uno por
\textit{(iii)}), se lo levanta a un haz invertible sobre $\mathfrak X$ y,
utilizando el teorema 1, se demuestra que un múltiplo de este último define una
inmersión de $\mathfrak X$ en la completación formal de un esquema
$\mathbf P_A^r$ («proyectivo tipo» de dimensión $r$ sobre $A$).

Para la demostración de los teoremas 1 a 3, véase [1].

\src{199}
\fsection{4}{Aplicación al teorema de conexión y al «Main theorem» de
Zariski}

Sea $f:X\longrightarrow Y$ un morfismo propio de esquemas. Por el teorema de
finitud, $f_*(\Osh_X)=\mathcal A$ es un haz coherente sobre $Y$; además, es un
haz de álgebras conmutativas, por lo que corresponde a un $Y$-esquema
$g:Y'\longrightarrow Y$ finito sobre $Y$ (definido por la condición de ser afín
sobre $Y$, es decir, que la imagen inversa de un abierto afín sea afín, y por
$g_*(\Osh_{Y'})=\mathcal A$). Es inmediato que $f$ se factoriza entonces
canónicamente como $f=gf'$, donde $f':X\longrightarrow Y'$ es un morfismo de
$X$ en $Y'$ que esta vez satisface $f'_*(\Osh_X)=\Osh_{Y'}$. Esta factorización
de $f$ se llama \emph{factorización de Stein} de $f$. Aplicando el primer
teorema de comparación, corolario 1, a $f'$ y a la parte de $Y'$ reducida a un
punto $y'$, se obtiene que $f'^{-1}(y')=X'$ es conexo (de otro modo, las
secciones formales de $X$ a lo largo de $X'$ no formarían un anillo local,
¡mientras que la completación de
$f'_*(\Osh_X)_{y'}=\Osh_{Y',y'}$ es local!). Hemos demostrado:

\result{TEOREMA}{5 («Teorema de conexión» de Zariski)} Sea
$f:X\longrightarrow Y$ un morfismo propio. Entonces $f$ se factoriza de manera
única (salvo isomorfismo) como $f=gf'$, donde $g:Y'\longrightarrow Y$ es finito
y $f':X\longrightarrow Y'$ satisface $f'_*(\Osh_X)=\Osh_{Y'}$ (de donde
$g_*(\Osh_{Y'})=f_*(\Osh_X)$). Las fibras de $f'$ son conexas; es decir, el
conjunto de las componentes conexas de una fibra $f^{-1}(y)$ de $f$ está en
correspondencia biunívoca con el conjunto de los puntos de $Y'$ sobre $y$, esto
es, con el conjunto de los ideales maximales de $f_*(\Osh_X)_y$.

De aquí se deducen inmediatamente las variantes habituales del teorema de
conexión. Señalemos aquí únicamente el

\src{200}
\result{COROLARIO}{1} Para que un punto $x$ de $X$ esté aislado en su fibra
$f^{-1}(y)$, es necesario y suficiente que la fibra $f'^{-1}(y')$ (donde
$y'=f'(x)$) se reduzca a $x$, o bien que $f'$ induzca un isomorfismo de un
entorno de $x$ sobre un entorno de $y'$. El conjunto de estos puntos es un
abierto $U$, y $f'$ induce un isomorfismo de $U$ sobre un abierto de $Y'$.

Para demostrar que $f'$ es un isomorfismo local en $x$, se observa que $f'$
induce un isomorfismo $\Osh_{y'}\longrightarrow\Osh_x$, como se ve gracias a
$f'_*(\Osh_X)=\Osh_{Y'}$ y al hecho de que (siendo $f'$ una aplicación cerrada
cuya fibra en $y'$ se reduce a $x$) los $f'^{-1}(V)$ recorren un sistema
fundamental de entornos de $x$ cuando $V$ recorre un sistema fundamental de
entornos de $y'$. De aquí se deduce enseguida el resultado siguiente, debido a
Chevalley en el caso «geométrico».

\result{COROLARIO}{2} Para que $f$ sea un morfismo finito, es necesario y
suficiente que sea propio y que sus fibras sean finitas.

Si así ocurre, $f'$ es efectivamente un isomorfismo por lo anterior.

Sea $f:X\longrightarrow Y$ un morfismo no necesariamente propio, pero
supongamos que $X$ está contenido como abierto en un $Y$-esquema propio
$\widetilde f:\widetilde X\longrightarrow Y$ (lo cual ocurre, en particular,
si $f$ es cuasiproyectivo). Aplicando el corolario 1, se ve que $\widetilde f'$
induce un isomorfismo del conjunto $U$ de los puntos de $X$ aislados en su fibra
sobre un abierto de $Y'$ (y que $U$ es, en efecto, abierto). Se deduce la
siguiente variante global del «Main Theorem» de Zariski.

\result{TEOREMA}{6 («Main Theorem» de Zariski)} Sea
$f:X\longrightarrow Y$ un morfismo de tipo finito. Entonces el conjunto $U$ de
los puntos de $X$ que están aislados en su fibra es abierto y, si $f$ es
separado,\footnote{La fe de erratas de 1962 sustituye aquí la hipótesis impresa
«$f$ cuasiproyectivo» por la hipótesis más débil «$f$ separado», en virtud de
SGA VIII, 6.2: todo morfismo cuasifinito y separado es cuasiproyectivo.} $U$ es
$Y$-isomorfo a una parte abierta de un esquema $Y'$ finito sobre $Y$.

Como un morfismo de tipo finito es localmente afín y, \textit{a fortiori},
localmente cuasiproyectivo, del teorema 6 se deducen enseguida las variantes
locales usuales del Main Theorem.

\fsection{5}{Aplicación al estudio cohomológico de los morfismos propios y
planos}

Sea $f:X\longrightarrow Y$ un morfismo propio, y sea $F$ un haz coherente sobre
$X$; se supone que $F$ es $Y$-plano, es decir, los $F_x$ son módulos planos sobre los
anillos $\Osh_y$ ($y=f(x)$). Esto significa también que, para todo $y\in Y$, si
filtramos $F$ a lo largo de la fibra $f^{-1}(y)$ mediante los
$\mathfrak m_y^nF$ (siendo $\mathfrak m_y$ el ideal maximal de $\Osh_y$), el
graduado asociado es isomorfo a
$(F/\mathfrak m_yF)\otimes_{\kappa(y)}\operatorname{gr}(\Osh_y)$; en otros
términos, se tiene
\[
  \mathfrak m_y^nF/\mathfrak m_y^{n+1}F
  =F_y\otimes_{\kappa(y)}
  (\mathfrak m_y^n/\mathfrak m_y^{n+1})
\]
para todo entero $n$, donde $X_y$ denota la fibra $f^{-1}(y)$ considerada como
esquema propio sobre el cuerpo residual $\kappa(y)$ de $y$, y $F_y$ el haz
$F/\mathfrak m_yF$ inducido por $F$ sobre $X_y$. Teniendo en cuenta este
isomorfismo y el teorema 2, se obtienen cotas y, en ocasiones, cálculos de los
$R^qf_*(F)$ en un entorno de $y$, conociendo la cohomología de $X_y$ con
coeficientes en $F_y$. El teorema 2 adopta aquí la forma
\[
  \widehat{R^qf_*(F)_y}
  =\varprojlim_n H^q(X_y,F/\mathfrak m_y^nF).
\]
Señalaremos únicamente la consecuencia siguiente:

\src{201}
\result{PROPOSICIÓN}{1} Sean $f:X\longrightarrow Y$ un morfismo propio y $F$
un haz sobre $X$, coherente y $Y$-plano. Sea $y\in Y$, sea $q$ un entero, y
supongamos que $H^q(X_y,F_y)=0$. Entonces $R^qf_*(F)$ es nulo en un entorno de
$y$ y, para todo $n$, el homomorfismo natural
\[
  R^{q-1}f_*(F)_y\longrightarrow
  H^{q-1}(X_y,F/\mathfrak m_y^nF)
\]
es sobreyectivo.

En particular, si $f$ es un morfismo plano (es decir, $\Osh_X$ es $Y$-plano),
entonces todo haz coherente localmente libre $F$ sobre $X$ es $Y$-plano. Sean
$F,G$ dos haces de este tipo; aplicando la proposición 1 a
$\Homsh_{\Osh_X}(F,G)$ y a $q=1$, se obtiene:

\result{TEOREMA}{7} Sea $f$ un morfismo propio y plano, sean $F,G$ dos haces
coherentes localmente libres sobre $X$, sea $y\in Y$ y supongamos que
\[
  H^1\bigl(X_y,\Homsh_{\Osh_{X_y}}(F_y,G_y)\bigr)=0,
\]
entonces todo homomorfismo $u_0:F_y\longrightarrow G_y$ está inducido por un
homomorfismo
\[
  u:F|V\longrightarrow G|V,
  \qquad V=f^{-1}(U),
\]
donde $U$ es un entorno de $y$.

\result{COROLARIO}{} Si $u_0$ es un isomorfismo (resp. un monomorfismo, un
epimorfismo), entonces lo mismo ocurre con $u$ para $U$ suficientemente
pequeño.

En particular:

\src{202}
\result{COROLARIO}{2} Sea $E_0$ un haz coherente localmente libre sobre $X_y$
tal que
\[
  H^1\bigl(X_y;\Homsh_{\Osh_{X_y}}(E_0,E_0)\bigr)=0,
\]
entonces dos haces localmente libres cuyas restricciones a $X_y$ sean isomorfas
a $E_0$ son isomorfos en un entorno de $X_y$.

Así:

\result{COROLARIO}{3} Supongamos $H^1(X_y,\Osh_{X_y})=0$. Entonces dos haces
invertibles sobre $X$ (es decir, localmente isomorfos a $\Osh_X$) cuyas
restricciones a $X_y$ sean isomorfas son isomorfos en un entorno de $X_y$.\footnote{El texto impreso no precisa el entorno. Se restituye aquí el alcance local del corolario 2, utilizado también en la demostración de la proposición 2.}

Se deduce:

\result{PROPOSICIÓN}{2} Sean $Y$ un esquema conexo y $E$ un haz coherente localmente libre de rango al menos $2$ sobre $Y$; consideremos el fibrado de espacios proyectivos
$X=\mathbf P(E)$ asociado a $E$, provisto de su conocido haz invertible
$\Osh_X(1)$. Entonces todo haz invertible $L$ sobre $X$ es isomorfo a un haz de
la forma
\[
  f^*(L')\otimes\Osh_X(n),
\]
donde $L'$ es un haz invertible sobre $Y$ y $n$ un entero. Este último está
determinado de manera única, y $L'$ está determinado salvo isomorfismo.

El corolario 3 precedente demuestra que $L$ es isomorfo a algún $\Osh_X(n)$ en
un entorno de cada fibra. El resto es casi formal.

La proposición 2 permite determinar los $Y$-morfismos de $X=\mathbf P(E)$ en
otro fibrado proyectivo. En particular, se obtiene:

\result{COROLARIO}{} Sea $u$ un automorfismo de $X=\mathbf P(E)$. Entonces
existen un haz invertible $L'$ sobre $Y$ y un isomorfismo $v$ de $E$ sobre
$E\otimes L'$ tales que $u$ es el isomorfismo correspondiente
\[
  \mathbf P(E)\xrightarrow{\sim}\mathbf P(E\otimes L')=\mathbf P(E);
\]
el par $(v,L')$ está determinado salvo isomorfismo.

Sea $\Gamma$ el conjunto de las clases de fibrados invertibles $L'$ sobre $Y$
tales que $E\otimes L'$ sea isomorfo a $E$. Sus elementos son de torsión, pues,
si $n$ es el rango de $E$, se ve (tomando las potencias exteriores $n$-ésimas)
que debe tenerse $L'^{\otimes n}\simeq\Osh_Y$. El corolario puede expresarse
entonces diciendo que se tiene una sucesión exacta de grupos
\[
  e\longrightarrow
  \operatorname{Aut}(E)/\Gamma(Y,\Osh_Y^*)
  \longrightarrow \operatorname{Aut}_Y(X)
  \longrightarrow \Gamma\longrightarrow e
\]
(deducida, por lo demás, de la sucesión exacta de cohomología obtenida a partir
de la siguiente sucesión exacta de haces de grupos
\[
  e\longrightarrow \Osh_Y^*
  \longrightarrow \underline{\operatorname{Aut}}(E)
  \longrightarrow \underline{\operatorname{Aut}}_Y(X)
  \longrightarrow e,
\]
donde $\Osh_Y^*$ es el haz de las «unidades» de $\Osh_Y$, identificado con el
centro de $\underline{\operatorname{Aut}}(E)$).

\src{203}
\fsection{6}{Aplicación a teoremas de existencia y unicidad para haces y
esquemas sobre un anillo completo para la topología $J$-ádica}

El teorema 7 daba un resultado de unicidad para haces coherentes localmente
libres utilizando los teoremas 1 y 2. Utilizando el teorema 3, deducimos ahora
teoremas de existencia de haces, de morfismos de esquemas o de esquemas. En lo
sucesivo, $A$ denota un anillo noetheriano local, separado y completo. El método
general consiste siempre en efectuar construcciones formales, lo cual consiste
esencialmente en hacer geometría algebraica sobre un anillo artiniano, y extraer
de ellas conclusiones de naturaleza «algebraica» mediante los tres teoremas
fundamentales.

\result{PROPOSICIÓN}{3} Sea $\mathfrak X$ un esquema formal propio y plano
sobre $A$, y sea $F_0$ un haz localmente libre sobre $X_0$ tal que
\[
  H^2\bigl(X_0,\Homsh_{\Osh_{X_0}}(F_0,F_0)\bigr)=0.
\]
Entonces existe un haz localmente libre $F$ sobre $\mathfrak X$ que induce
sobre $X_0$ un haz isomorfo a $F_0$. (Además, este $F$ es único salvo
isomorfismo si
$H^1(X_0,\Homsh_{\Osh_{X_0}}(F_0,F_0))=0$.)

Se construyen sucesivamente haces localmente libres $F_n$ sobre los $X_n$, cada
uno de los cuales induce el anterior. La construcción de $F_n$ encuentra una
obstrucción en
\[
  H^2\bigl(X_0,\Homsh_{\Osh_{X_0}}(F_0,F_0)\bigr)
  \otimes_{A/J}(J^n/J^{n+1}),
\]
que es, por tanto, nula por hipótesis. Utilizando ahora el teorema 3, se
obtiene:

\result{COROLARIO}{1} Sea $X$ un esquema propio y plano sobre $A$, y sea $F_0$
como arriba. Entonces existe un haz localmente libre $F$ sobre $X$ que induce
sobre $X_0$ un haz isomorfo a $F_0$. Además, este $F$ es único salvo
isomorfismo si
$H^1(X_0,\Homsh_{\Osh_{X_0}}(F_0,F_0))=0$.

Sea $X_0$ un esquema de tipo finito sobre el cuerpo $k$; se supone que $X_0$ es
simple (por lo cual entendemos absolutamente simple) sobre $k$, pero no
necesariamente propio sobre $k$. Sea $A$ un anillo artiniano local de cuerpo
residual $k$. Nos interesa hallar los esquemas $X$ planos sobre $A$ tales que
$X\otimes_A k=X_0$ (este es el punto de partida de la «teoría de módulos» o de
las «variaciones de estructura» de $X_0$). Es equivalente dar un tal $X$, o dar
sobre el espacio topológico $X_0$ un haz $\Osh_X$ provisto de las estructuras
siguientes:
\begin{enumerate}
  \item[\textit{i.}] Es un haz de $A$-álgebras.
  \item[\textit{ii.}] Está provisto de un homomorfismo de aumentación
    $\Osh_X\longrightarrow\Osh_{X_0}$ (compatible con las estructuras de
    $A$-álgebras); estos datos están sujetos a las condiciones siguientes: la
    aumentación induce un isomorfismo
    $\Osh_X\otimes_A k\xrightarrow{\sim}\Osh_{X_0}$; $\Osh_X$ es plano sobre
    $A$, es decir, el graduado asociado a $\Osh_X$, filtrado por las potencias
    del ideal maximal $\mathfrak m$ de $A$, es isomorfo a
    $\operatorname{gr}^0(\Osh_X)\otimes_k\operatorname{gr}(A)$; es decir, se
    tienen isomorfismos
    \[
      \mathfrak m^n\Osh_X/\mathfrak m^{n+1}\Osh_X
      =\Osh_{X_0}\otimes_k(\mathfrak m^n/\mathfrak m^{n+1}).
    \]
\end{enumerate}
El hecho fundamental es el siguiente:

\src{204}
\result{TEOREMA}{8} Sea $X_0$ un esquema de tipo finito y simple sobre el
cuerpo $k$, y supongamos que $X_0$ es afín. Sea $A$ un anillo local artiniano de
cuerpo residual $k$. Entonces existe un $A$-esquema $X$ plano sobre $A$ tal que
$X\otimes_A k=X_0$, y dos esquemas de este tipo son necesariamente isomorfos.

Obsérvese que el isomorfismo en cuestión no es canónico, pues en general $X$
tendrá $A$-automorfismos no triviales que inducen la identidad sobre $X_0$. Por
otra parte, en general no hay una elección «canónica» de un $X$ que satisfaga
las condiciones dadas, salvo cuando $A$ es una $k$-álgebra (caso de
características iguales), en cuyo caso puede tomarse $X=X_0\otimes_k A$, es
decir, $\Osh_X=\Osh_{X_0}\otimes_k A$ (sea o no afín $X_0$). En el caso de
características distintas, ignoro en general, cuando $X_0$ no es afín, si se
puede «levantar» $X_0$ a un $X$ definido sobre $A$. Sin embargo, sea $n$ un
entero $>0$, pongamos $A_{n-1}=A/\mathfrak m^n$, y supongamos que se ha
levantado $X_0$ a un $A_{n-1}$-esquema plano $X_{n-1}$; se pretende levantar
$X_{n-1}$ a un $A_n$-esquema plano $X_n$. Por el teorema 8 ya se sabe que esto
es posible localmente; por otra parte, se comprueba fácilmente que, si $U_n$
levanta un abierto $U_{n-1}$ de $X_{n-1}$, entonces el haz de los grupos de
automorfismos de $U_n$ (que inducen la identidad sobre $U_{n-1}$) es
canónicamente isomorfo a
\[
  \Theta_{X_0/k}\otimes_k(\mathfrak m^n/\mathfrak m^{n+1})
\]
restringido a $U_{n-1}$ y, en particular, es conmutativo
($\Theta_{X_0/k}$ denota el haz de gérmenes de $k$-derivaciones sobre $X_0$).
Se sigue fácilmente que hay una obstrucción para construir $X_n$ que levante
$X_{n-1}$, la cual se encuentra en
\[
  H^2(X_0,\Theta_{X_0/k})\otimes_k
  \mathfrak m^n/\mathfrak m^{n+1}.
\]
Por consiguiente:

\result{COROLARIO}{1} Sea $X_0$ un esquema de tipo finito simple sobre $k$, y
supongamos $H^2(X_0,\Theta_{X_0/k})=0$. Entonces, para todo anillo local
artiniano $A$ de cuerpo residual $k$, existe un $A$-esquema plano $X$ tal que
$X\otimes_A k=X_0$.

Además, si se ha podido hallar un $X$ plano sobre $A$ que levante $X_0$,
entonces, en virtud del teorema 8, el conjunto de las clases (salvo isomorfismo)
de $A$-esquemas planos que levantan $X_0$ se identifica con
\[
  H^1\bigl(X_0,\underline{\operatorname{Aut}}(X)\bigr),
\]
donde, naturalmente, $\underline{\operatorname{Aut}}(X)$ denota el haz de
gérmenes de automorfismos del haz de $A$-álgebras $\Osh_X$ compatibles con la
aumentación. La filtración de $\Osh_X$ define una filtración de
$\underline{\operatorname{Aut}}(X)$, y el cociente de este haz por el $n$-ésimo
subgrupo de la filtración es $\underline{\operatorname{Aut}}(X_n)$; el graduado
asociado a esta filtración es conmutativo y se identifica con
$\Theta_{X_0/k}\otimes_k\operatorname{gr}(A)$. En particular, si
$\mathfrak m^{n+1}$ es la primera potencia de $\mathfrak m$ que es nula,
entonces $F^n(\underline{\operatorname{Aut}}(X))$ (el último término de la
filtración) pertenece al centro de $\underline{\operatorname{Aut}}(X)$ y es
isomorfo a $\Theta_{X_0/k}\otimes_k\mathfrak m^n$; es también el haz de
gérmenes de automorfismos de $X$ que inducen la identidad sobre
$X_{n-1}=X\otimes_A A/\mathfrak m^n$. Utilizando estos resultados se obtienen
enseguida los enunciados siguientes:

\src{205}
\result{COROLARIO}{2} Sea $X_0$ un esquema de tipo finito simple sobre $k$, sea
$A$ un anillo local artiniano de cuerpo residual $k$ e ideal maximal
$\mathfrak m$, supongamos $\mathfrak m^{n+1}=0$, sea
$A_{n-1}=A/\mathfrak m^n$ y, por último, sea $X_{n-1}$ un
$A_{n-1}$-esquema plano tal que $X_{n-1}\otimes_A k=X_0$. Entonces el conjunto
de las clases (salvo un isomorfismo que induzca la identidad sobre $X_{n-1}$)
de $A$-esquemas planos $X_n$ tales que
$X_n\otimes_A A_{n-1}=X_{n-1}$ es vacío o es un espacio homogéneo principal
bajo
\[
  H^1(X_0,\Theta_{X_0/k})\otimes_k\mathfrak m^n.
\]

(Obsérvese que, en general, no hay un elemento origen privilegiado en este
último espacio, pues no hay una manera privilegiada de levantar $X_{n-1}$ a
$X_n$).

\result{COROLARIO}{3} Sea $X_0$ un esquema de tipo finito simple sobre $k$, y
supongamos $H^1(X_0,\Theta_{X_0/k})=0$. Entonces, para todo anillo local
artiniano $A$ de cuerpo residual $k$, existe a lo sumo (salvo isomorfismo) un
$A$-esquema plano $X$ tal que $X\otimes_A k=X_0$.

Los corolarios 1 y 3 implican inmediatamente los enunciados, en apariencia más
generales, que se obtienen suponiendo solamente que $A$ es un anillo local
noetheriano completo de cuerpo residual $k$, con tal de introducir allí $X$
como un esquema formal sobre $A$:

\result{TEOREMA}{9} Sean $k$ un cuerpo y $X_0$ un esquema de tipo finito simple
sobre $k$. Para todo anillo local noetheriano completo $A$ de cuerpo residual
$k$, sea $F(A)$ el conjunto de las clases (salvo un isomorfismo que induzca la
identidad sobre $X_0$) de esquemas formales $\mathfrak X$ sobre $A$, de tipo
finito y planos sobre $A$, tales que $\mathfrak X\otimes_A k=X_0$. Con estas
notaciones:
\begin{enumerate}
  \item[\textit{i.}] Si $H^1(X_0,\Theta_{X_0/k})=0$, entonces, para todo $A$,
    $F(A)$ tiene a lo sumo un elemento.
  \item[\textit{ii.}] Si $H^2(X_0,\Theta_{X_0/k})=0$, entonces, para todo $A$,
    $F(A)$ tiene al menos un elemento.
\end{enumerate}

\src{206}
\result{COROLARIO}{1} Supongamos que $X_0$ es propio sobre $k$. Bajo la
condición \textit{(i)}, para todo $A$ existe a lo sumo (salvo un isomorfismo que
induzca la identidad sobre $X_0$) un esquema $X$ propio y plano sobre $A$ tal
que $X\otimes_A k=X_0$.

Se utiliza el teorema 3, corolario 2. Por ejemplo:

\result{COROLARIO}{2} Si $X$ es un esquema propio y plano sobre $A$ tal que
$X\otimes_A k$ es isomorfo al esquema proyectivo tipo $\mathbf P_k^r$ de
dimensión $r$ sobre $k$, entonces $X$ es isomorfo a $\mathbf P_A^r$.

(También se puede deducir este resultado de la proposición 3, corolario 1).

\result{COROLARIO}{3} Sea $X_0$ un esquema simple y proyectivo sobre $k$, y
supongamos que se tiene
\[
  H^2(X_0,\Osh_{X_0})=H^2(X_0,\Theta_{X_0/k})=0.
\]
Entonces, para todo $A$, existe un esquema proyectivo y plano $X$ sobre $A$ tal
que $X\otimes_A k=X_0$.

Se combinan el teorema 9\textit{(ii)} y el teorema 4. En particular:

\result{COROLARIO}{4} Sea $X_0$ el esquema de una curva algebraica completa y
simple sobre $k$. Entonces, para todo anillo local noetheriano completo $A$ de
cuerpo residual $k$, existe un «esquema de curva simple» $X$ sobre $A$ tal que
$X\otimes_A k=X_0$.

\result{OBSERVACIONES}{}

1) Los corolarios 3 y 4 son especialmente interesantes si $k$ es de
característica $p\ne0$, tomando como $A$ un anillo de valoración discreta de
característica $0$ y cuerpo residual $k$; por ejemplo, el «menor $A$ posible»,
a saber, aquel cuyo ideal maximal está generado por $p$. (De hecho, por los
teoremas de Cohen, basta conocer los corolarios 3 y 4 para un anillo $A$ de este
tipo). Señalemos a este respecto que, según los especialistas, no se sabe si
existen esquemas sobre un cuerpo $k$ que no sean la reducción módulo $p$ de un
esquema plano definido sobre un tal anillo $A$. Al menos, los resultados de este
número proporcionan un medio de investigar sistemáticamente esta cuestión. Se
debería comenzar por examinar si la primera obstrucción que aparece en
$H^2(X_0,\Theta_{X_0/k})$ es necesariamente nula.

\emph{Complemento de 1962.} --- J.-P. Serre ha construido una variedad
proyectiva no singular, de dimensión 3, sobre un cuerpo algebraicamente cerrado
$k$ de característica $p>0$, que no se obtiene por reducción de un esquema
propio sobre un anillo local íntegro de cuerpo residual $k$ y cuerpo de
fracciones de característica cero. Mumford habría encontrado un ejemplo análogo
con una superficie proyectiva no singular.

2) Obsérvese que el teorema 3, y la técnica correspondiente, solo es válido para
un anillo de base completo (local, para fijar ideas).

Para pasar de resultados conocidos para la completación de un anillo local a
los resultados correspondientes para ese mismo anillo local, se necesitaría un
cuarto «teorema fundamental», cuyo enunciado definitivo aún está por hallar.

\src{207}
3) Compárense los resultados de este número (en particular, los corolarios 1 y
2 precedentes) y del siguiente con los resultados de Kodaira--Spencer sobre la
variación de las estructuras complejas. Utilizando el teorema conjetural al que
se acaba de aludir, bajo las condiciones del corolario 1, pero sin suponer ya
que $A$ sea completo, se debería poder concluir que existe un anillo $A'$ que
contiene a $A$, finito y no ramificado sobre $A$, tal que $X\otimes_A A'$ y
$X'\otimes_A A'$ son $A'$-isomorfos (siendo $X,X'$ dos $A$-esquemas propios y
planos dados, tales que $X\otimes_A k=X'\otimes_A k=X_0$). Esto es al menos lo
que puede demostrarse cuando $X_0=\mathbf P_k^r$, utilizando el corolario de la
proposición 2. En todo caso, cuando $H^1(X_0,\Theta_{X_0/k})=0$, puede
demostrarse que las fibras de $X,X'$ sobre todo punto $y$ de
$Y=\operatorname{Spec}(A)$ son isomorfas, al menos después de pasar a la
clausura algebraica del cuerpo residual $\kappa(y)$. (Se tiene un resultado
local, aparentemente más fuerte, sin suponer ya necesariamente que $A$ sea
local). En cuanto a las «variaciones de estructura» del espacio proyectivo,
señalemos además la cuestión siguiente, sugerida por un problema correspondiente
de Kodaira--Spencer. Sea $X$ un esquema propio y plano sobre el anillo local
íntegro $A$ de cuerpo de fracciones $K$ y cuerpo residual $k$, y supongamos que
$X\otimes_A K$ es isomorfo a $\mathbf P_K^r$; ¿es cierto que
$X\otimes_A k=X_0$ es isomorfo a $\mathbf P_k^r$ (al menos sobre la clausura
algebraica de $k$)? En este problema se puede suponer que $A$ es un anillo de
valoración discreta completo. Se plantea un problema análogo cuando $X_0$ es
una variedad abeliana.

\emph{Complemento de 1962.} --- El conferenciante es ahora menos optimista
respecto de los resultados conjeturados arriba. No obstante, Hironaka ha
resuelto afirmativamente el problema relativo a las variaciones de estructura
del espacio proyectivo, y Koizumi el problema análogo para las variedades
abelianas.

\fsection{7}{Aplicación a la «Teoría de módulos»}

Como el propio conferenciante apenas acaba de abordar esta teoría, nos veremos
obligados a limitarnos a algunas indicaciones. Para simplificar, nos situamos
sobre un cuerpo $k$, es decir, trabajamos en igual característica, aunque el
teorema 8 permite tratar también el caso general sin cambio esencial, al
parecer. Actualmente no hemos superado la etapa «formal»; sin embargo, el
conferenciante espera poder construir a partir de ella verdaderos esquemas de
módulos en ciertos casos y, en particular, construir para cada entero $g$ un
esquema sobre los enteros que desempeñe el papel de esquema de módulos universal
para las curvas simples de género $g$.

\emph{Complemento de 1962.} --- Mumford ha construido los esquemas de módulos
para las curvas de género $g$ (cf. SHMI). El teorema 10 demuestra además que
los «esquemas de Jacobi de nivel $n$» de la teoría de Módulos son no singulares,
e incluso simples sobre $\mathbf Z$.

\src{208}
Retomemos la situación y las notaciones del teorema 9, suponiendo ahora que $A$
es un álgebra local de rango finito sobre $k$, que para simplificar se supone
algebraicamente cerrado. Entonces $F(A)$ puede considerarse como un funtor
covariante de $A$, con valores en la categoría de conjuntos: un homomorfismo de
$k$-álgebras $A\longrightarrow B$ define una aplicación
$F(A)\longrightarrow F(B)$, pues todo $A$-esquema plano $X$ tal que
$X\otimes_A k=X_0$ da lugar a un $B$-esquema $X\otimes_A B$ con las mismas
propiedades. Supongamos que se puede hallar una $k$-álgebra local noetheriana
completa $\mathcal O$ y un isomorfismo funtorial
\begin{equation}
  \tag{*}
  \operatorname{Hom}(\mathcal O,A)\xrightarrow{\sim}F(A)
\end{equation}
(donde el primer miembro denota los homomorfismos de $k$-álgebras). Se ve
fácilmente que un tal $\mathcal O$ está determinado salvo isomorfismo canónico;
se llama entonces al espectro formal $\mathfrak Y$ de $\mathcal O$ (es decir,
al espacio topológico reducido a un punto y provisto de un haz de anillos
topológicos reducido a $\mathcal O$) el \emph{esquema formal de módulos} para
$X_0$. (Obsérvese que no existe necesariamente). Sea $\mathfrak r$ el ideal
maximal de $\mathcal O$; para cada $n$, sea
$\mathcal O_n=\mathcal O/\mathfrak r^{n+1}$ (por tanto $\mathcal O_0=k$).
Entonces el homomorfismo canónico $\mathcal O\longrightarrow\mathcal O_n$ es un
elemento de $\operatorname{Hom}(\mathcal O,\mathcal O_n)$, por lo que define un
elemento de $F(\mathcal O_n)$, es decir, un $\mathcal O_n$-esquema plano $X_n$
cuya reducción módulo $\mathfrak r$ es $X_0$. Estos $X_n$ se deducen unos de
otros por extensión de escalares (es decir, aquí por reducciones), de donde
resulta que provienen de un esquema formal $\mathfrak X$ bien determinado sobre
el esquema formal de módulos $\mathfrak Y$; $\mathfrak X$ es plano sobre
$\mathfrak Y$ y $\mathfrak X_0=X_0$. Como se ve enseguida, el isomorfismo (*)
se obtiene entonces asociando a todo homomorfismo de $k$-álgebras
$\mathcal O\longrightarrow A$ la clase del $A$-esquema
$\mathfrak X\otimes_{\mathcal O}A$ (es decir, a todo morfismo de $k$-esquemas
$Y'=\operatorname{Spec}(A)\longrightarrow\mathfrak Y$ se asocia el
$Y'$-esquema $\mathfrak X\otimes_{\mathfrak Y}Y'$ obtenido por cambio de base).
Además, se ve que el isomorfismo (*) y su descripción precedente siguen siendo
válidos si solo se supone que $A$ es una $k$-álgebra local noetheriana y
completa (no necesariamente artiniana). Naturalmente, como de costumbre,
$\mathcal O$ puede tener perfectamente \textit{a priori} elementos nilpotentes,
y parece probable que existan casos en que $\mathcal O$ sea él mismo artiniano
sin ser idéntico a $k$. Esto muestra hasta qué punto el punto de vista de
Kodaira--Spencer (que se limita a tomar como $A$ anillos regulares) es
inadecuado \textit{a priori} en el caso general.

\src{209}
Queda por dar condiciones suficientes para que exista un esquema formal de
módulos para $X_0$, que se supone propio sobre $k$. En general, es fácil dar
condiciones necesarias y suficientes sencillas sobre un funtor
$A\longmapsto F(A)$ (de $k$-álgebras locales de rango finito, con valores en
los conjuntos) para que pueda escribirse en la forma
$\operatorname{Hom}(\mathcal O,A)$ para un $\mathcal O$ adecuado. No las
detallaremos aquí. Señalemos solamente que, en el caso que nos ocupa, estas
condiciones imponen sobre $X_0$ condiciones cohomológicas no triviales, y parece
poco probable que se satisfagan siempre, aunque el conferenciante no ha
construido ningún contraejemplo. En cambio, parece plausible que la condición
$H^0(X_0,\Theta_{X_0/k})=0$ sea una condición suficiente (aunque en modo alguno
necesaria) para la existencia de un esquema formal de módulos. Nos limitamos a
enunciar aquí un teorema en un caso particularmente sencillo (cuyo análogo en
la teoría de espacios analíticos es bien conocido, cf. Kodaira--Spencer), que se
establece sin dificultad mediante los resultados del número anterior:

\result{TEOREMA}{10} Sea $X_0$ un esquema propio y simple sobre el cuerpo $k$
tal que
\[
  H^0(X_0,\Theta_{X_0/k})=H^2(X_0,\Theta_{X_0/k})=0.
\]
Entonces existe un esquema formal de módulos para $X_0$, correspondiente a un
anillo local regular $\mathcal O$ (es decir, un álgebra de series formales sobre
$k$).

Como ya hemos señalado, no es cierto en general que el esquema formal
$\mathfrak X$ sobre $\mathcal O$ sea algebraizable; pero se sabe que sí lo es
cuando $X_0$ es proyectivo y $H^2(X_0,\Osh_{X_0})=0$ (teorema 4), por ejemplo
cuando $X_0$ tiene dimensión 1. Esto da cierto fundamento a la esperanza de
construir un esquema de módulos sobre los enteros para las curvas de género
dado\ldots{}

Obsérvese también que métodos como los expuestos en este número pueden aplicarse
a la construcción y el estudio de las variedades de Picard, así como a muchas
otras construcciones. Volveremos sobre ello próximamente.

\fsection{8}{Aplicación al grupo fundamental}

Las técnicas expuestas permiten abordar el estudio sistemático del grupo
fundamental siguiendo el modelo de la teoría topológica. Los dos primeros
teoremas enunciados en este número son generalizaciones de resultados de un
trabajo reciente de Lang--Serre.

Sea $X$ un esquema. Un $X$-esquema $X'$ se llama revestimiento no ramificado de
$X$ si
\begin{enumerate}
  \item[\textit{i.}] $X'$ es finito sobre $X$, es decir, está definido por un
    haz coherente de álgebras $\mathcal A=\mathcal A(X')$ sobre $X$;
  \item[\textit{ii.}] $\mathcal A$ es un haz localmente libre sobre $X$;
  \item[\textit{iii.}] Para todo $x\in X$, el cociente
    \[
      \mathcal A_x/\mathfrak m_x\mathcal A_x
      =\mathcal A_x\otimes_{\Osh_x}\kappa(x)
    \]
    es un álgebra separable sobre $\kappa(x)$.
\end{enumerate}

\src{210}
Esta noción de revestimiento no ramificado (debida a Serre y al conferenciante)
posee todas las propiedades elementales que cabe razonablemente esperar, cuya
lista no daremos. Limitémonos a decir que da lugar a una teoría de Galois
modelada sobre la teoría de Galois clásica (y que la contiene; las
demostraciones son más bien más sencillas que las comúnmente aceptadas para
esta última) y sobre la teoría de Galois de los revestimientos topológicos. Más
precisamente, llamemos \emph{punto geométrico} de un esquema $X$ a un morfismo
$a$ del espectro $\xi$ de un cuerpo algebraicamente cerrado $\Omega$ en $X$, es
decir, al dato de una extensión algebraicamente cerrada del cuerpo residual
$\kappa(x)$ de un punto $x=|a|$ de $X$ (llamado \emph{localidad} del punto
geométrico $a$). Si $X'$ es entonces un revestimiento no ramificado de $X$, se
le puede asociar el conjunto $E_a(X')$ de los «puntos geométricos de $X'$ sobre
$a$», es decir, el conjunto de los pares formados por un $x'\in X'$ sobre $x$
y un $\kappa(x)$-homomorfismo $\kappa(x')\longrightarrow\Omega$. Se obtiene así
(para $(X,a)$ fijos) un funtor $F(X,a)$ de la categoría $R(X)$ de los
revestimientos no ramificados $X'$ de $X$ en la categoría de los conjuntos
finitos. Si $X$ es conexo, el par formado por $R(X)$ y $F(X,a)$ posee las
propiedades formales necesarias para ser isomorfo al par análogo definido por
un grupo topológico compacto totalmente discontinuo $\pi$ adecuado (es decir,
un límite proyectivo de grupos finitos): se toma la categoría $C(\pi)$ de los
conjuntos finitos $E$ sobre los cuales $\pi$ actúa continuamente, y el funtor
identidad $F(\pi)(E)=E$ de esta categoría en la categoría de los conjuntos
finitos. Además, el grupo $\pi$ está determinado salvo isomorfismo canónico por
la condición de que $(C(\pi),F(\pi))$ sea isomorfo a un par dado. En este caso,
$\pi$ se llama \emph{grupo fundamental} del esquema conexo $X$ en el punto
geométrico $a$, y se denota por $\pi_1(X,a)$. Si $X$ no es conexo, se lo
reemplaza por la componente conexa de $x=|a|$. Si, por el contrario, $X$ es
conexo, los grupos $\pi_1(X,a)$ correspondientes a dos puntos geométricos
$a',a''$ de $X$ son isomorfos (el isomorfismo está determinado salvo
automorfismos interiores), de modo que, como de costumbre, puede escogerse $a$
de la manera más conveniente, por ejemplo en el punto genérico de $X$, supuesto
irreducible. Naturalmente, $\pi_1(X,a)$ es un funtor covariante del esquema
punteado $(X,a)$. Todo enunciado relativo a la clasificación de los
revestimientos no ramificados\footnote{El texto impreso dice «revestimientos
inseparables»; el párrafo trata de la categoría $R(X)$ de los revestimientos no
ramificados y de su equivalencia con los conjuntos finitos dotados de una acción
continua de $\pi_1(X,a)$.} puede traducirse entonces al lenguaje de la teoría de
grupos siguiendo el diccionario bien conocido (salvo que hay que tener en cuenta
que aquí se trata de grupos topológicos).

\src{211}
Nuestro objetivo es desarrollar el análogo de la sucesión exacta de homotopía
de los espacios fibrados para un morfismo propio $f:X\longrightarrow Y$.
Evidentemente, al no saber qué son los grupos de homotopía superiores, solo
obtendremos resultados necesariamente incompletos. Para poder aplicar los
teoremas fundamentales del n\textsuperscript{o} 3, debemos explicitar primero
algunos lemas elementales relativos a los esquemas sobre cuerpos o anillos de
Artin (¡conforme al procedimiento general!).

\result{LEMA}{1} Sea $(X',a')$ un revestimiento no ramificado punteado asociado
a una representación punteada de $\pi_1(X,a)$ en un conjunto finito $E$
(provisto del punto marcado $e$). Entonces el morfismo canónico
$\pi_1(X',a')\longrightarrow\pi_1(X,a)$ identifica el primer grupo con el
estabilizador de $e$ en $\pi_1(X,a)$ (y es, por tanto, inyectivo).

\result{LEMA}{2} Sea $X$ un esquema algebraico sobre el cuerpo $k$, y sea $k'$
una extensión radicial de $k$. Entonces todo revestimiento no ramificado de
$X\otimes_k k'$ proviene por imagen inversa (es decir, por extensión de
escalares) de un revestimiento no ramificado de $X$, determinado salvo
isomorfismo canónico.

De estos dos lemas resulta, en particular, que para toda extensión algebraica
$K$ de $k$, y para todo punto geométrico $a'$ de $X'=X\otimes_k K$ que se
proyecta sobre el punto geométrico $a$ de $X$, el homomorfismo funtorial
$\pi_1(X',a')\longrightarrow\pi_1(X,a)$ es inyectivo.

\result{LEMA}{3} Sea $X$ un esquema completo sobre un anillo local artiniano
$A$, tal que $H^0(X,\Osh_X)=A$; sea $X'$ un revestimiento no ramificado de $X$,
y sea $A'=H^0(X',\Osh_{X'})$, que es, por tanto, un anillo finito sobre $A$
(que \textit{a priori} puede estar ramificado sobre $A$). Sean $X_0,X'_0$ los
subesquemas reducidos asociados a $X,X'$ (obtenidos dividiendo por los haces de
elementos nilpotentes de $\Osh_X,\Osh_{X'}$), y sea $k$ un subcuerpo de
$A/\mathfrak r(A)$ sobre el cual $A/\mathfrak r(A)$ sea finito (así, $X_0$ es
un esquema algebraico completo sobre $k$, y $X'_0$ es un revestimiento no
ramificado suyo). Sea, por último, $\Omega$ una extensión algebraicamente
cerrada de $k$, y consideremos el revestimiento no ramificado
$X'_0\otimes_k\Omega$ de $X_0\otimes_k\Omega$. Entonces las dos condiciones
siguientes son equivalentes:
\begin{enumerate}
  \item[\textit{i.}] $X'_0\otimes_k\Omega$ está completamente descompuesto
    sobre $X_0\otimes_k\Omega$.
  \item[\textit{ii.}] El morfismo natural
    $X'\longrightarrow X\otimes_A A'$ es un isomorfismo.
\end{enumerate}
Bajo estas condiciones, $A'$ es una extensión no ramificada de $A$. Por último,
si $X'$ es conexo, la condición \textit{(i)} equivale a la condición siguiente,
en apariencia más débil:
\begin{enumerate}
  \item[\textit{i bis.}] $X'_0\otimes_k\Omega$ admite una sección regular
    sobre $X_0\otimes_k\Omega$.
\end{enumerate}
Cuando se satisface la condición \textit{(ii)}, diremos que el revestimiento no
ramificado $X'$ de $X$ es \emph{geométricamente trivial}.

\src{212}
\result{LEMA}{4} Sea $f:X\longrightarrow Y$ un morfismo propio tal que
$f_*(\Osh_X)=\Osh_Y$. Sean $a$ un punto geométrico de $X$ y $b$ su proyección
sobre $Y$. Entonces $\pi_1(X,a)\longrightarrow\pi_1(Y,b)$ es sobreyectivo.

En efecto, hay que demostrar lo siguiente: si un revestimiento no ramificado
$Y'$ de $Y$ (correspondiente a un haz de álgebras localmente libre
$\mathcal A$) es tal que $X\otimes_Y Y'$ es disconexo, entonces $Y'$ también lo
es. En efecto, $f^*(\mathcal A)$ será suma directa de dos haces de anillos no
nulos y, por tanto, también lo será su imagen directa; ahora bien, esta última
no es otra que $\mathcal A\otimes f_*(\Osh_X)=\mathcal A$.

\result{LEMA}{5} Sea $X$ un esquema completo sobre un cuerpo $k$; se supone que
$H^0(X,\Osh_X)$ es un anillo local $A$ y que $A/\mathfrak r(A)$ es radicial
sobre $k$. Sea $\Omega$ una clausura algebraica de $k$, y sea
$\overline X=X\otimes_k\Omega$ (es conexo). Elijamos un punto geométrico
$\overline a$ de $\overline X$ que se proyecte sobre el punto geométrico $a$ de
$X$; entonces se tiene una sucesión exacta
\[
  e\longrightarrow\pi_1(\overline X,\overline a)
  \longrightarrow\pi_1(X,a)
  \longrightarrow\pi_1(k,b)\longrightarrow e
\]
(donde $\pi_1(k,b)$ es el grupo de Galois de $\Omega$ sobre $k$).

La inyectividad del primer homomorfismo se ha visto con los lemas 1 y 2, la
exactitud en el medio resulta del lema 3 y, por último, la sobreyectividad del
último homomorfismo (que es la única que utiliza el hecho de que
$A/\mathfrak r(A)$ sea radicial) resulta del lema 4.

\result{PROPOSICIÓN}{4} Sea $f:X\longrightarrow Y$ un morfismo propio y plano
tal que, para todo $y\in Y$, $H^0(f^{-1}(y),\Osh_{f^{-1}(y)})$ sea un álgebra
separable sobre el cuerpo residual $\kappa(y)$ (lo cual ocurre, por ejemplo, si
$f^{-1}(y)$ es un esquema separable sobre $\kappa(y)$, es decir, reducido y tal
que los cuerpos correspondientes a sus componentes irreducibles sean
extensiones separables de $\kappa(y)$). Entonces el revestimiento $Y'$ de $Y$
asociado a $f_*(\Osh_X)$ es no ramificado.

La demostración es fácil gracias al teorema 2.

\src{213}
Esta proposición, junto con el lema 1, reduce prácticamente el estudio
homotópico de los morfismos propios y planos (con fibras separables) al caso en
que $f_*(\Osh_X)=\Osh_Y$ (pues, utilizando la factorización de Stein, se
reemplaza $Y$ por $Y'$).

\result{OBSERVACIÓN}{} Un morfismo plano de tipo finito cuyas fibras son
esquemas separables (resp. simples) se llama \emph{separable} (resp.
\emph{simple}). Se demuestra que, si $f$ es plano y $f^{-1}(y)$ es separable
(resp. simple), existe un entorno de $f^{-1}(y)$ en el que $f$ es separable
(resp. simple). El mismo resultado es válido para «absolutamente normal»
(«teorema de Bertini»).

Sea $f:X\longrightarrow Y$ un morfismo propio tal que
\begin{equation}
  \tag{i}
  f_*(\Osh_X)=\Osh_Y,
\end{equation}
y sea $X'$ un esquema finito sobre $X$. Sea $Y'$ el revestimiento de $Y$
correspondiente a la factorización de Stein de $X'\longrightarrow Y$
(cf. teorema 5). Sea $y\in Y$; entonces el conjunto de las componentes conexas
de la fibra $F'$ de $X'$ sobre $y$ se identifica con el conjunto de los puntos
$y'$ de $Y'$ sobre $y$ (teorema 5). Consideremos el morfismo evidente
\begin{equation}
  \tag{*}
  X'\longrightarrow X\times_Y Y'
\end{equation}
deducido de los morfismos naturales $X'\longrightarrow X$ y
$X'\longrightarrow Y'$\footnote{El texto impreso omite la prima en el segundo
$X'$; el morfismo hacia $Y'$ es el de la factorización de Stein de
$X'\to Y$.}; será un isomorfismo siempre que $X'$ sea de la forma
$X\times_Y Y''$, donde $Y''$ es un revestimiento no ramificado de $Y$, y
entonces $Y'$ no será sino $Y''$ y (*) será la identidad. Queremos dar
precisamente condiciones bajo las cuales $X'$ sea de la forma que acabamos de
indicar, es decir, que $Y'$ sea no ramificado y (*) un isomorfismo. Para ello,
introduzcamos la fibra $F$ de $X$ en $y$; es un esquema propio sobre
$\kappa(y)$ del cual $F'$ es un revestimiento (no ramificado si $X'$ lo es).
Sea $F'_1$ una componente conexa de $F'$ correspondiente a un punto $y'_1$ de
$Y'$ sobre $y$. Supongamos además
\begin{enumerate}
  \item[\textit{(ii)}] $X'$ es no ramificado sobre $X$ en los puntos de $F'_1$
    (por tanto, $F'_1$ es un revestimiento no ramificado de $F$), y
  \item[\textit{(iii)}] $F'_1$ es un revestimiento geométricamente trivial de
    $F$ (cf. lema 3).
\end{enumerate}

\result{TEOREMA}{11} En estas condiciones, existe un entorno abierto $U'$ de
$y'_1$ en $Y'$ tal que (*) es un isomorfismo sobre $U'$. Además, $Y'$ es no
ramificado en $y'_1$ sobre $Y$ (pero puede ser ramificado en los demás puntos
$y'$ de $Y'$ sobre $y$).

Naturalmente, las condiciones \textit{(ii)} y \textit{(iii)} son también
necesarias para que sea válida la conclusión del teorema. La demostración del
teorema es fácil, con ayuda del lema 3 y del teorema 2.

\src{214}
\result{COROLARIO}{1} Supongamos siempre satisfecha \textit{(i)}. Para que un
revestimiento no ramificado sobre $X$ sea isomorfo a la imagen inversa de un
revestimiento no ramificado $Y'$ de $Y$, es necesario y suficiente que $X'$
induzca sobre cada fibra $f^{-1}(y)$ un revestimiento geométricamente trivial.

Como, según el teorema 11, el conjunto de los puntos de $Y$ para los cuales se
cumple esta condición es abierto, basta verificarla en los puntos $y$ que son
cerrados\ldots{} Señalemos el enunciado siguiente, equivalente al corolario 1:
el núcleo del homomorfismo (sobreyectivo según el lema 4)
$\pi_1(X)\longrightarrow\pi_1(Y)$ es el subgrupo invariante cerrado engendrado
por las imágenes en $\pi_1(X)$ de los $\pi_1(\overline{f^{-1}(y)})$, donde
\[
  \overline{f^{-1}(y)}
  =f^{-1}(y)\otimes_{\kappa(y)}\overline{\kappa(y)},
\]
$\overline{\kappa(y)}$ designa una clausura algebraica de $\kappa(y)$. Se
observará que, al no poder elegir el mismo punto base para todas las fibras,
los homomorfismos $\pi_1(\overline{f^{-1}(y)})\longrightarrow\pi_1(X)$ solo
quedan determinados en cualquier caso (una vez elegido un punto base para $X$
y, por tanto, para $Y$) módulo composición por un automorfismo interior en
$\pi_1(X)$.

\result{COROLARIO}{2} En las condiciones generales del teorema 11, supongamos
además que $Y,X,X'$ son íntegros, y sean $K,L,L'$ sus cuerpos. Entonces existe
una subextensión separable $K'$ de $K$ en $L'$, linealmente disjunta de $L$, tal
que $L'=LK'$ (de donde $L'=L\otimes_K K'$).

(Se aplica la última parte del lema 3 a la fibra geométrica de $X$). El caso de
aplicación más interesante del teorema 11 se obtiene cuando $f$ es un morfismo
separable. Entonces $X'$ es también separable sobre $Y$; por tanto, en virtud
de la proposición 4, $Y'$ es no ramificado sobre $Y$, y el segundo miembro
$X\times_Y Y'$ de (*) es no ramificado sobre $X$. Se concluye fácilmente:

\result{COROLARIO}{3} Supongamos, además de \textit{(i)}, que $f$ es separable.
Sea $X'$ un revestimiento no ramificado y conexo de $X$; para que
$X'$\footnote{El texto impreso dice aquí «$X$»; el sujeto del criterio es el
revestimiento $X'$ introducido en la frase anterior.} sea imagen inversa de un
revestimiento no ramificado $Y'$ de $Y$, es necesario y suficiente que el
revestimiento $\overline F'$ inducido sobre una fibra geométrica
$\overline F=f^{-1}(y)$ admita una sección regular.

Se observará que no era necesario suponer que $\overline F'$ fuera
geométricamente trivial sobre $\overline F$ (lo cual será cierto \textit{a
posteriori}, aunque \textit{a priori} esta condición sea mucho más fuerte). El
corolario 3 es, de hecho, equivalente al enunciado siguiente:

\src{215}
\result{COROLARIO}{4} Sea $f:X\longrightarrow Y$ un morfismo propio y separable
tal que $f_*(\Osh_X)=\Osh_Y$, sea $\overline F$ la fibra geométrica de un punto
$y\in Y$, y elijamos un punto geométrico en $\overline F$, del cual se obtienen,
mediante los morfismos $\overline F\longrightarrow X\longrightarrow Y$, puntos
geométricos en $X,Y$, que se tomarán como puntos base para los grupos
fundamentales de $\overline F,X,Y$. En estas condiciones se tiene la sucesión
exacta
\[
  \boxed{
    \pi_1(\overline F)\longrightarrow\pi_1(X)
    \longrightarrow\pi_1(Y)\longrightarrow e
  }.
\]

De ello se deducen fácilmente los dos enunciados siguientes de Serre--Lang,
liberados aquí de toda hipótesis de normalidad:

\result{COROLARIO}{5} Sean $X,Y$ esquemas conexos sobre un cuerpo $k$;
supongamos que el esquema reducido $X_{\mathrm{réd}}$ es separable sobre $k$
(lo cual es automáticamente cierto si $k$ es perfecto), y que $X$ o $Y$ es
propio sobre $k$.\footnote{Condición añadida en la fe de erratas de 1962.}
Elijamos un punto geométrico $a$ (resp. $b$) en $X$ (resp. $Y$), de donde se
obtienen un punto geométrico $c=(a,b)$ en $X\times_kY$ y un morfismo natural
\[
  \pi_1(X\times_kY,c)\longrightarrow
  \pi_1(X,a)\times\pi_1(Y,b)
\]
(deducido de los morfismos funtoriales de $\pi_1(X\times Y,c)$ en $\pi_1(X,a)$
y $\pi_1(Y,b)$). Este morfismo es inyectivo, e incluso biyectivo si $k$ es
algebraicamente cerrado.

(La sobreyectividad en este último caso es casi trivial). De ello se deduce,
con Serre--Lang:

\result{COROLARIO}{6} Sea $X$ un esquema algebraico conexo y propio sobre un cuerpo
algebraicamente cerrado $k$, y sea $K$ una extensión algebraicamente cerrada
de $k$; entonces los grupos fundamentales de $X$ y $X\otimes_kK$ son los
mismos, es decir, todo revestimiento no ramificado de este último esquema
proviene por extensión de escalares de un revestimiento no ramificado (único
salvo isomorfismo) de $X$.

\result{OBSERVACIONES}{}

1) Utilizando la proposición 4, se ve que la hipótesis
$f_*(\Osh_X)=\Osh_Y$ del corolario 4 no es esencial. En el caso general, en vez
de colocar el grupo unidad $e$ después de $\pi_1(Y)$, hay que continuar con
\[
  \pi_0(\overline F)\longrightarrow\pi_0(X)
  \longrightarrow\pi_0(Y)\longrightarrow e.
\]
como en Topología algebraica.

\src{216}
2) En general, por el momento no se puede decir nada sobre el núcleo de
$\pi_1(\overline F)\longrightarrow\pi_1(X)$, en el que debería intervenir un
$\pi_2(Y)$. Parece, sin embargo, que debería poder demostrarse que
$\pi_1(\overline F)\longrightarrow\pi_1(X)$ es inyectivo si $Y$ es el espectro
de un anillo local $A$, apoyándose en el teorema 12 de más abajo (que afirma que
así es si $A$ es completo).

El teorema 11 solo utilizaba los teoremas 1 y 2. Vamos a aplicar ahora el
teorema 3, apoyándonos en el lema elemental siguiente:

\result{LEMA}{6} Sea $X$ un esquema, y $X_0$ el esquema reducido correspondiente
(es decir, en el que se han anulado los elementos nilpotentes). Entonces todo
revestimiento no ramificado $X'_0$ de $X_0$ está inducido por un revestimiento
no ramificado $X'$ de $X$, determinado salvo isomorfismo canónico.

Este lema, de naturaleza puramente local, desempeña aquí un papel análogo al del
teorema 8 en teoría de módulos. Combinándolo con el teorema de existencia
(teorema 3), se deduce:

\result{TEOREMA}{12} Sea $A$ un anillo local noetheriano y completo de cuerpo
residual $k$. Sea $X$ un esquema propio sobre $A$. Entonces todo revestimiento
no ramificado $X'_0$ de $X_0=X\otimes_A k$ está inducido por un revestimiento no
ramificado $X'$ de $X$, único salvo isomorfismo.

En otros términos:

\result{COROLARIO}{1} Elijamos un punto geométrico en $X_0$ como punto base para
los grupos fundamentales de $X_0$ y $X$. Entonces el homomorfismo canónico
$\pi_1(X_0)\longrightarrow\pi_1(X)$ es un isomorfismo.

Apliquemos ahora a $X_0$ el lema 5 (suponiendo, para simplificar, que
$H^0(X_0,\Osh_{X_0})=k$), y observemos que, por ser $A$ completo, sus
extensiones no ramificadas corresponden a las extensiones no ramificadas de su
cuerpo residual, es decir, se tiene $\pi_1(Y)=\pi_1(k)$ (donde
$Y=\operatorname{Spec}(A)$). Se obtiene la sucesión exacta:
\[
  e\longrightarrow\pi_1(\overline X_0)
  \longrightarrow\pi_1(X)\longrightarrow\pi_1(Y)\longrightarrow e.
\]

\result{COROLARIO}{2} Sea $f:X\longrightarrow Y$ un morfismo propio y plano,
sean $y_1$ un punto de $Y$ e $y_0$ una especialización de $y_1$, y consideremos
las fibras «geométricas» correspondientes $\overline X_1$ y $\overline X_0$;
se supone que $\overline X_0$ es separable y conexa (lo que implica que
$\overline X_1$ satisface las mismas condiciones). Entonces se puede hallar un
homomorfismo de grupos
$\pi_1(\overline X_1)\longrightarrow\pi_1(\overline X_0)$, definido salvo un
automorfismo interior, y este homomorfismo es sobreyectivo.

\src{217}
Podría esperarse que este homomorfismo fuera siempre biyectivo; por desgracia,
en general no es así si $\kappa(y_0)$ tiene característica $>0$. Obtendremos,
sin embargo, más abajo una cota para el núcleo de este homomorfismo (al menos si
$\overline X_0$ es simple), de la que se sigue que, si $\kappa(y_0)$ tiene
característica $0$, el homomorfismo anterior es biyectivo (resultado que también
podría demostrarse por vía trascendente). En todo caso, obtenemos ya una cota
para el $\pi_1$ de una fibra especial por medio del de una fibra genérica.
Utilizando, por ejemplo, el hecho de que una curva algebraica en característica
$p$ se puede levantar a una curva en característica $0$ (teorema 9, corolario
4), se obtiene por vía trascendente:

\result{COROLARIO}{3} Sea $X_0$ el esquema de una curva completa simple sobre
un cuerpo algebraicamente cerrado de característica cualquiera, y sea $g$ el
género de $X_0$; entonces $\pi_1(X_0)$ admite $2g$ generadores topológicos,
ligados por la relación bien conocida.

De ello se deduce, mediante una técnica bien conocida de secciones
hiperplanas:

\result{COROLARIO}{4} Sea $X$ un esquema proyectivo simple sobre un cuerpo
algebraicamente cerrado de característica cualquiera; entonces $\pi_1(X)$
admite un número finito de generadores topológicos.

Busquemos el núcleo del homomorfismo
$\pi_1(\overline X_1)\longrightarrow\pi_1(\overline X_0)$. Para ello se puede
suponer que $Y$ es el espectro de un anillo de valoración discreta completo
$V=\Osh_y$ ($y=y_0$). La cuestión equivale a la siguiente: dado un revestimiento
no ramificado $\overline X'_1$ de $\overline X_1$ (que, si se quiere, puede
suponerse galoisiano), determinar en qué condiciones proviene de un
revestimiento no ramificado de $X_0$. \textit{A priori}, el revestimiento dado
proviene, por extensión de escalares, de un revestimiento no ramificado $X'_1$
de $X_1\otimes_K K'$, donde $K'$ es una subextensión finita de la clausura
algebraica $\overline K$ del cuerpo de fracciones $K$ de $V$; si
$\overline X'_1$ fuera galoisiano de grupo $G$, se puede elegir $X'_1$
galoisiano de grupo $G$. Se tiene entonces: para que
$X'_1\otimes_{K'}\overline K=\overline X'_1$ provenga de un revestimiento no
ramificado de $X_0$, es necesario y suficiente que se pueda hallar una extensión
finita $K''$ de $K'$ en $\overline K$ tal que
$X''_1=X'_1\otimes_{K'}K''$ sea de la forma $X''\otimes_{V''}K''$, donde $V''$
es la clausura normal de $V$ en $K''$ y $X''$ es un revestimiento no ramificado
de $X\otimes_VV''$. Supongamos, por ejemplo, que $X_0$ es absolutamente normal;
de ello resulta que $X\otimes_VV''$ es normal (pues es plano sobre $V''$ y tiene
fibra especial normal), y su cuerpo de funciones es idéntico al $K''(X_1)$ de
\[
  X_1\otimes_KK''
  =(X\otimes_VK)\otimes_KK''
  =X\otimes_VK''
  =(X\otimes_VV'')\otimes_{V''}K''.
\]
Sea $L''=K''(X'_1)$ el cuerpo de funciones de $X'_1\otimes_{K'}K''$; es una
extensión finita separable de $K''(X_1)$, y la condición anterior significa
también que $L''$ es una extensión no ramificada del cuerpo de funciones de
$X\otimes_VV''$ (es decir, la normalización de $X\otimes_VV''$ en $L''$ es no
ramificada sobre $X\otimes_VV''$). Basta, por lo demás, verificar que $L''$ es
no ramificado en los puntos de la fibra especial de $X\otimes_VV''$ (pues es no
ramificado sobre la fibra genérica $X_1\otimes_KK''$). Si ahora $X_0$ es simple,
del «teorema de pureza» de Nagata--Zariski resulta incluso que basta verificar
que $L''$ es no ramificado sobre el anillo local $O''$ del punto genérico de la
fibra especial de $X\otimes_VV''$, que es un anillo de valoración discreta,
normalización en $K''(X_1)$ del anillo local $O\subset K(X_1)$ del punto
genérico de la fibra especial de $X$. Por tanto, se ha obtenido:

\src{218}
\result{COROLARIO}{5} En las condiciones y con las notaciones anteriores, para
que el revestimiento no ramificado
$X'_1\otimes_{K'}\overline K$ de
$\overline X_1=X_1\otimes_K\overline K$ provenga de un revestimiento no
ramificado de $\overline X_0$, es necesario y suficiente que exista una
subextensión finita $K''$ de $\overline K/K'$ tal que $K''(X'_1)$ sea no
ramificado sobre el anillo de valoración discreta $O''\subset K''(X_1)$.

Observemos ahora que $O''$ es la normalización en $K''(X_1)$ del anillo de
valoración discreta $O'\subset K'(X_1)$ (normalización de $O$ en $K'(X_1)$), y
que $O'$ contiene la normalización $V'$ de $V$ en $K'$; un uniformizante $u$ de
$V'$ es también un uniformizante de $O'$. Supongamos ahora que $X'_1$ es
galoisiano de grupo de Galois $G$, de orden $n$ primo con la característica $p$
de $\kappa(y_0)$ (que es también la característica del cuerpo residual de
$O'$). Por tanto, $K'(X'_1)$ está «moderadamente ramificado» sobre $O'$, de lo
que resulta fácilmente («lema de Abhyankar») que, si se le adjunta una raíz
$n$-ésima $v$ de un uniformizante de $O'$, se vuelve no ramificado sobre la
normalización de $O'$ en $K'(X_1)(v)$. Ahora bien, se puede tomar como $v$ una
raíz $n$-ésima de un uniformizante de $V'$, lo que prueba que se satisface la
condición del corolario 5. (Esta posibilidad de utilizar el lema de Abhyankar y
el teorema de pureza me la ha vendido Serre). Para expresar el resultado
obtenido, introduzcamos, para todo grupo compacto totalmente discontinuo $\pi$,
el grupo cociente $\widetilde\pi$ de $\pi$ por el subgrupo cerrado engendrado
por sus $p$-subgrupos de Sylow, es decir, el límite proyectivo de los grupos
cocientes discretos de $\pi$ que son de orden primo con $p$. Con esta notación,
se obtiene:

\result{TEOREMA}{13} Sea $f:X\longrightarrow Y$ un morfismo propio y plano,
sean $y_1$ un punto de $Y$ e $y_0$ una especialización de $y_1$, y supongamos
$\overline X_0$ conexo y simple. Entonces el homomorfismo
$\widetilde\pi_1(\overline X_1)\longrightarrow
\widetilde\pi_1(\overline X_0)$ deducido del homomorfismo sobreyectivo del
teorema 12, corolario 2, es un isomorfismo.

En otros términos:

\src{219}
\result{COROLARIO}{1} La clasificación de los revestimientos galoisianos no
ramificados, de grupo de Galois de orden primo con la característica $p$ de
$\kappa(y_0)$, es la misma para $\overline X_0$ y para $\overline X_1$.

En particular, si $\kappa(y_0)$ tiene característica cero, se obtiene por vía
algebraica que $\pi_1(\overline X_1)\longrightarrow\pi_1(\overline X_0)$ es
biyectivo.

Señalemos, por último, que las técnicas utilizadas proporcionan también el
resultado siguiente, más general que el teorema 13:

\result{TEOREMA}{14} Sea $f:X\longrightarrow Y$ un morfismo propio y simple,
sea $D$ un subesquema cerrado de $X$, simple sobre $Y$, de codimensión 1 en
todos sus puntos. Para una fibra $Z=f^{-1}(z)$ de $f$, sea $Z'=Z-Z\cap D$ y sea
$\pi_1^t(Z')$ el cociente del grupo fundamental $\pi_1(Z')$ que clasifica los
revestimientos no ramificados de $Z'$ que están «moderadamente ramificados»
sobre $\overline{Z\cap D}$. Sean $y_0,y_1$ como en el teorema 13. Entonces se
tiene un homomorfismo sobreyectivo (definido salvo un automorfismo interior)
\[
  \pi_1^t(\overline X'_1)\longrightarrow\pi_1^t(\overline X'_0),
\]
y el homomorfismo correspondiente
\[
  \widetilde\pi_1^t(\overline X'_1)
  \longrightarrow\widetilde\pi_1^t(\overline X'_0)
\]
es un isomorfismo.

De ello se deducen las variantes correspondientes de los corolarios del
teorema 13 y del corolario 4 del teorema 12. Del mismo modo, utilizando el
teorema 9, corolario 3, se obtiene por vía trascendente:

\result{COROLARIO}{} Sea $X_0$ el esquema de una curva completa simple sobre un
cuerpo algebraicamente cerrado de característica cualquiera, y sea
$S=(s_i)_{1\le i\le n}$ una parte finita con $n$ elementos de $X_0$. Entonces
$\pi_1^t(X_0-S)$ admite $2g+n$ generadores topológicos
$x_i,y_i$ ($1\le i\le g$) y $\sigma_j$ ($1\le j\le n$), ligados por la relación
\[
  \left(\prod_i x_i y_i x_i^{-1}y_i^{-1}\right)
  \sigma_1\cdots\sigma_n=1,
\]
donde los $\sigma_j$ son generadores de los grupos de inercia correspondientes
a los $s_j$. Para todo grupo finito $G$ de orden primo con la característica,
engendrado por elementos $x_i,y_i,\sigma_j$ que satisfacen la relación anterior,
existe un revestimiento galoisiano no ramificado de $X_0-S$, de grupo $G$, cuyos
grupos de inercia en los puntos $s_j$ están engendrados por los $\sigma_j$.

Cuando $X_0$ es de género 0 y $n=3$, se tiene una solución del «problema de los
tres puntos», al menos para los revestimientos galoisianos de orden primo con
la característica. (Aquí, por lo demás, el teorema 9 es inútil, y parece, por
otra parte, que el corolario anterior puede deducirse del caso particular
considerado en el problema de los tres puntos).

\src{220}
\result{OBSERVACIONES}{}

1) Un estudio más completo, en el que sin duda intervendrían revestimientos
galoisianos generalizados de $X,X_0,X_1$ (con grupo de Galois eventualmente
infinitesimal), debería permitir recuperar el núcleo del teorema 12, corolario
2. En cambio, el estudio de los revestimientos que admiten ramificación no
«moderada» parece tener que ser mucho más difícil.

2) El lema 6, junto con un resultado de Grauert relativo a la completación
formal de un esquema proyectivo no singular a lo largo de una sección
hiperplana (o con el teorema, no demostrado por el momento, mencionado en la
observación 2 después del teorema 11), permite también demostrar en geometría
algebraica «abstracta» el teorema clásico de Lefschetz sobre el grupo
fundamental.

\begin{center}
  \textbf{BIBLIOGRAFÍA}
\end{center}

\begin{description}
  \item[{[1]}] Dieudonné (J.) y Grothendieck (A.). --- \emph{Éléments de
    géométrie algébrique}, Publications mathématiques de l'Institut des Hautes
    Études Scientifiques (de próxima aparición).
  \item[{[2]}] Grothendieck (Alexander). --- \emph{The cohomology theory of
    abstract algebraic varieties}. International Congress of Mathematicians
    [1958. Edinburgh] (de próxima aparición).
  \item[{[3]}] Serre (Jean-Pierre). --- \emph{Faisceaux algébriques
    cohérents}, Annals of Math., Series 2, t. 61, 1955, p. 197--278.
  \item[{[4]}] Serre (Jean-Pierre). --- \emph{Géométrie algébrique et
    géométrie analytique}, Ann. Institut Fourier Grenoble, t. 6, 1955--1956,
    p. 1--42.
\end{description}
